\section*{Capítulo \ref{ch:b2:multint} --- Integrais de linha e integrais múltiplas}

\begin{solution}{exo:b2:multint:1}
(a) Segmento $\gamma(t) = (t, t)$, $t \in [0,1]$:
\[
\int_\gamma y^2\dd x + x\dd y
= \int_0^1 (t^2 + t)\,\dd t = \frac13 + \frac12 = \frac56 .
\]
(b) Parábola $\gamma(t) = (t, t^2)$:
\[
\int_0^1 \bigl(t^4\cdot 1 + t\cdot 2t\bigr)\dd t
= \frac15 + \frac23 = \frac{13}{15} .
\]
Os dois valores diferem, de modo que a integral depende do caminho: a forma
\emph{não} é \omterm{def:b2:multint:exact}{exata} --- coerentemente, $P_y = 2y \neq 1 = Q_x$, logo
ela não é sequer fechada.
\end{solution}

\begin{solution}{exo:b2:multint:2}
$P = 2xy + y^3$, $Q = x^2 + 3xy^2 + 1$: $P_y = 2x + 3y^2 = Q_x$,
fechada em $\R^2$. Procure $f$ com $f_x = P$: $f = x^2y + xy^3
+ g(y)$; então $f_y = x^2 + 3xy^2 + g'(y) = Q$ força $g'(y) = 1$,
digamos $g(y) = y$. Assim
\[
f(x, y) = x^2y + xy^3 + y
\]
é um \omterm{def:b2:multint:exact}{potencial} ($\R^2$ é estrelado, de modo que um \omterm{def:b2:multint:exact}{potencial} tinha de
existir pelo lema de Poincaré --- mas exibi-lo é mais rápido). Pelo
\cref{thm:b2:multint:ftc}, para qualquer arco de $(0,0)$ a $(1,2)$:
\[
\int_\gamma\omega = f(1, 2) - f(0, 0) = 2 + 8 + 2 = 12 .
\]
\end{solution}

\begin{solution}{exo:b2:multint:3}
O domínio é $\{0 \leq x \leq 1,\ x^2 \leq y \leq x\}$.
Primeiro em $y$:
\[
\int_0^1\!\int_{x^2}^{x}(x + y)\,\dd y\,\dd x
= \int_0^1\Bigl(x(x - x^2)
+ \frac{x^2 - x^4}{2}\Bigr)\dd x
= \int_0^1\Bigl(\frac{3x^2}{2} - x^3 - \frac{x^4}{2}\Bigr)\dd x
= \frac12 - \frac14 - \frac1{10} = \frac{3}{20} .
\]
Primeiro em $x$: a fatia à altura $y \in [0, 1]$ é $y \leq x \leq
\sqrt y$, logo
\[
\int_0^1\!\int_{y}^{\sqrt y}(x + y)\,\dd x\,\dd y
= \int_0^1\Bigl(\frac{y - y^2}{2}
+ y(\sqrt y - y)\Bigr)\dd y
= \frac14 - \frac16 + \frac25 - \frac13 = \frac{3}{20} .
\]
\end{solution}

\begin{solution}{exo:b2:multint:4}
\emph{Primeira integral.} No disco $D_R$, em coordenadas polares:
\[
\iint_{D_R}\frac{\dd x\,\dd y}{(1 + x^2 + y^2)^2}
= \int_0^{2\pi}\!\!\int_0^R
\frac{\rho\,\dd\rho\,\dd\alpha}{(1 + \rho^2)^2}
= 2\pi\Bigl[-\frac{1}{2(1 + \rho^2)}\Bigr]_0^R
= \pi\Bigl(1 - \frac{1}{1 + R^2}\Bigr)
\xrightarrow[R \to \infty]{} \pi .
\]
\emph{Segunda integral.} O quarto de disco é $0 \leq \alpha \leq
\frac\pi2$, $0 \leq \rho \leq 1$, e $xy =
\rho^2\cos\alpha\sin\alpha$:
\[
\iint_{D'}xy\,\dd x\,\dd y
= \int_0^{\pi/2}\!\!\cos\alpha\sin\alpha\,\dd\alpha
\int_0^1 \rho^3\,\dd\rho
= \frac12\cdot\frac14 = \frac18 .
\]
\end{solution}

\begin{solution}{exo:b2:multint:5}
Pelo~\cref{cor:b2:multint:area} com $x = \cos^3 t$, $y = \sin^3 t$:
$x' = -3\cos^2 t\sin t$, $y' = 3\sin^2 t\cos t$, logo
\[
xy' - yx' = 3\cos^4 t\sin^2 t + 3\sin^4 t\cos^2 t
= 3\sin^2 t\cos^2 t = \frac{3}{4}\sin^2 2t
= \frac{3}{8}(1 - \cos 4t) .
\]
Portanto
\[
A = \frac12\int_0^{2\pi}\frac38(1 - \cos 4t)\,\dd t
= \frac{3}{16}\cdot 2\pi = \frac{3\pi}{8} .
\]
(O \omterm{pb:b2:curves:1}{astroide} cabe no disco unitário de \omterm{def:b2:multint:domain}{área} $\pi$; três oitavos de
$\pi$ é plausível para sua forma de estrela de quatro cúspides.)
\end{solution}

\begin{solution}{exo:b2:multint:6}
\emph{Empilhamento:} acima de cada $(x, y)$ do disco unitário $D$, $z$
vai de $x^2 + y^2$ a $1$:
\[
V = \iint_D \bigl(1 - x^2 - y^2\bigr)\dd x\,\dd y
= \int_0^{2\pi}\!\!\int_0^1 (1 - \rho^2)\rho\,\dd\rho\,\dd\alpha
= 2\pi\Bigl(\frac12 - \frac14\Bigr) = \frac\pi2 .
\]
\emph{Fatiamento:} a fatia à altura $z \in [0, 1]$ é o disco
$x^2 + y^2 \leq z$, de \omterm{def:b2:multint:domain}{área} $\pi z$:
\[
V = \int_0^1 \pi z\,\dd z = \frac\pi2 .
\]
\end{solution}

\begin{solution}{exo:b2:multint:7}
Em coordenadas esféricas o elemento de volume é
$r^2\cos\varphi\,\dd r\,\dd\theta\,\dd\varphi$
(\cref{ex:b2:multint:spherical}), e a parte angular integra
a $4\pi$ ($2\pi$ vindo de $\theta$, $\int_{-\pi/2}^{\pi/2}\cos = 2$).
O volume da casca é
\[
\int_a^b 4\pi r^2\,\dd r = \frac43\pi\bigl(b^3 - a^3\bigr).
\]
Para a segunda integral, o integrando $1/r$ depende apenas de $r$:
\[
\iiint_B \frac{\dd x\,\dd y\,\dd z}{r}
= \int_0^R 4\pi r^2\cdot\frac1r\,\dd r
= 4\pi\,\frac{R^2}{2} = 2\pi R^2 .
\]
(O integrando explode na origem, mas inofensivamente: $r^2/r = r$
é \omterm{def:b2:metric:continuity}{contínua} --- a integral sobre as cascas $\varepsilon \leq r
\leq R$ converge quando $\varepsilon \to 0$, que é o sentido preciso
do enunciado. Esse tipo de cálculo é o primeiro
passo rumo ao teorema de Newton segundo o qual uma bola homogênea atrai como
uma massa pontual em seu centro.)
\end{solution}

\begin{solution}{exo:b2:multint:8}
O integrando $g(t; x, y) = xP(tx, ty) + yQ(tx, ty)$ é
$\mathcal{C}^1$ em $(x, y)$, \omterm{def:b2:metric:continuity}{contínuo} em $t$, com derivadas
parciais \omterm{def:b2:metric:continuity}{contínuas} em $[0,1] \times U$; a derivação sob
o sinal de integral (\cref{ch:b2:integration}, aplicada no
intervalo \omterm{def:b2:metric:compact}{compacto} de $t$, $[0,1]$, sendo a dominação automática ali)
dá
\[
f_x(x, y)
= \int_0^1 \bigl(P(tx, ty) + tx\,P_x(tx, ty)
+ ty\,Q_x(tx, ty)\bigr)\dd t .
\]
Usando o fato de a forma ser fechada, $Q_x = P_y$:
\[
tx\,P_x(tx, ty) + ty\,P_y(tx, ty)
= t\,\frac{\dd}{\dd t}\bigl[P(tx, ty)\bigr] ,
\]
de modo que o integrando é $P(tx, ty) + t\frac{\dd}{\dd t}P(tx, ty) =
\frac{\dd}{\dd t}\bigl[t\,P(tx, ty)\bigr]$ e
\[
f_x(x, y) = \Bigl[t\,P(tx, ty)\Bigr]_0^1 = P(x, y) .
\]
Simetricamente, $f_y = Q$ (mesmo cálculo com $P_y = Q_x$ usado
no outro sentido). Note onde a hipótese entra: $f$ é definida
integrando ao longo do segmento $[0, M]$, que está em $U$
precisamente porque $U$ é estrelado.
\end{solution}

\begin{solution}{exo:b2:multint:9}
Na faixa $[0, A] \times [0, \infty)$ a função $(x, y)
\mapsto e^{-xy}\sin x$ não é \omterm{def:b2:series:def}{absolutamente} integrável até $y =
\infty$ \omterm{def:b2:funcseq:def}{uniformemente} num sentido ingênuo, mas cada integral iterada
converge e sua igualdade decorre de Fubini em $[0, A]
\times [0, B]$ mais um limite $B \to \infty$ (a cauda
$\int_0^A\int_B^\infty e^{-xy}\abs{\sin x}\,\dd y\,\dd x \leq
\int_0^A \frac{e^{-Bx}\abs{\sin x}}{x}\dd x \leq \int_0^A e^{-Bx}
\dd x\to 0$, usando $\abs{\sin x} \leq x$).

\emph{Primeiro em $y$:} $\int_0^\infty e^{-xy}\,\dd y = \frac1x$ para $x
> 0$, de modo que a primeira integral vale $\int_0^A \frac{\sin x}{x}\,\dd x$.

\emph{Primeiro em $x$:} duas integrações por partes (ou tomar a
parte imaginária de $\int_0^A e^{(i - y)x}\dd x$) dão
\[
\int_0^A e^{-xy}\sin x\,\dd x
= \frac{1 - e^{-Ay}(y\sin A + \cos A)}{1 + y^2} .
\]
Integrando em $y$ sobre $[0, \infty)$, o termo $\int_0^\infty
\frac{\dd y}{1 + y^2} = \frac\pi2$ se destaca:
\[
\int_0^A \frac{\sin x}{x}\,\dd x
= \frac{\pi}{2}
- \int_0^\infty e^{-Ay}\,\frac{y\sin A + \cos A}{1 + y^2}\,\dd y .
\]
O resto é majorado por $\int_0^\infty e^{-Ay}\frac{y +
1}{1 + y^2}\dd y \leq \int_0^\infty e^{-Ay}\cdot\frac{1+y}{1+y^2}
\,\dd y \to 0$ quando $A \to \infty$ (convergência dominada, ou a
majoração grosseira $\frac{1 + y}{1 + y^2} \leq \frac32$, que dá
$\frac{3}{2A}$). Portanto $\int_0^\infty\frac{\sin x}{x}\dd x =
\frac\pi2$ --- o mesmo valor obtido no
\cref{ch:b2:integration} derivando uma \omterm{thm:b2:integration:continuity}{integral com parâmetro};
aqui Fubini faz o trabalho no lugar.
\end{solution}

\begin{solution}{exo:b2:multint:10}
Parametrize por \omterm{def:b2:curves:length}{comprimento de arco} $s \in [0, 2\pi]$, de modo que $x'^2 + y'^2 = 1$,
e translade para que $\int_0^{2\pi} x(s)\,\dd s = 0$. Pelo
\cref{cor:b2:multint:area},
\[
2A = \oint x\,\dd y - y\,\dd x
= \int_0^{2\pi}\bigl(xy' - yx'\bigr)\dd s .
\]
Integrando $\oint y\,\dd x$ por partes sobre o período (os termos de
bordo se cancelam por periodicidade), $-\int yx' = \int y'x$, de modo que, de fato,
$2A = 2\int_0^{2\pi}xy'\,\dd s$. Então $2xy' \leq x^2 + y'^2$
dá
\[
2A \leq \int_0^{2\pi}\bigl(x^2 + y'^2\bigr)\dd s
= \int_0^{2\pi} x^2 + \int_0^{2\pi}\bigl(1 - x'^2\bigr)
= 2\pi - \int_0^{2\pi}\bigl(x'^2 - x^2\bigr)\dd s .
\]
A desigualdade de Wirtinger (exercícios do~\cref{ch:b2:fourier}: para uma
função $\mathcal{C}^1$ de período $2\pi$ e média nula,
$\int x^2 \leq \int x'^2$) torna a última integral não negativa:
$A \leq \pi$. A igualdade exige igualdade em Wirtinger ($x(s) =
a\cos s + b\sin s$) e em $2xy' \leq x^2 + y'^2$ ($y' = x$
\omterm{def:b2:funcseq:def}{pontualmente}), o que força $y = a\sin s - b\cos s + c$: a curva é
o círculo unitário (convenientemente centrado). Como uma curva de \omterm{def:b2:curves:length}{comprimento} $L$
se reescala para \omterm{def:b2:curves:length}{comprimento} $2\pi$, o enunciado geral é $A \leq
\frac{L^2}{4\pi}$: entre todas as curvas fechadas de perímetro dado, o
círculo encerra a maior \omterm{def:b2:multint:domain}{área}.
\end{solution}

\begin{solution}{exo:b2:multint:11}
Em coordenadas esféricas $z = r\sin\varphi$ e $\dd x\,\dd
y\,\dd z = r^2\cos\varphi\,\dd r\,\dd\theta\,\dd\varphi$:
\[
\iiint_B z^2
= \int_0^R r^4\,\dd r\int_0^{2\pi}\dd\theta
\int_{-\pi/2}^{\pi/2}\sin^2\varphi\cos\varphi\,\dd\varphi
= \frac{R^5}5\cdot2\pi\cdot
\Bigl[\frac{\sin^3\varphi}3\Bigr]_{-\pi/2}^{\pi/2}
= \frac{4\pi R^5}{15}.
\]
Pela simetria da bola sob permutação das coordenadas,
$\iiint_B x^2 = \iiint_B y^2 = \iiint_B z^2$, logo $\iiint_B(x^2
+ y^2 + z^2) = 3\cdot\frac{4\pi R^5}{15} = \frac{4\pi R^5}5$.
Verificação por cascas: $\int_0^R r^2\cdot 4\pi r^2\,\dd r = \frac{4\pi
R^5}5$ --- o integrando $r^2$ é constante na esfera de
raio $r$, de \omterm{def:b2:multint:domain}{área} $4\pi r^2$.
\end{solution}

\begin{solution}{exo:b2:multint:12}
Parametrize a aresta de $(x_i, y_i)$ a $(x_{i+1}, y_{i+1})$
por $\gamma(t) = \bigl((1-t)x_i + tx_{i+1},\ (1-t)y_i +
ty_{i+1}\bigr)$. Sua contribuição a
$\frac12\oint(x\,\dd y - y\,\dd x)$ é
\[
\frac12\int_0^1\Bigl(\bigl((1-t)x_i +
tx_{i+1}\bigr)(y_{i+1} - y_i) - \bigl((1-t)y_i +
ty_{i+1}\bigr)(x_{i+1} - x_i)\Bigr)\dd t ,
\]
e, como $\int_0^1\bigl((1-t)u + tv\bigr)\dd t = \frac{u +
v}2$, isso vale
\[
\frac14\Bigl((x_i + x_{i+1})(y_{i+1} - y_i) - (y_i +
y_{i+1})(x_{i+1} - x_i)\Bigr)
= \frac12\bigl(x_iy_{i+1} - x_{i+1}y_i\bigr),
\]
cancelando-se os termos cruzados. Somando sobre as $m$ arestas obtém-se
a fórmula do cadarço, pelo~\cref{cor:b2:multint:area}. Triângulo
$(0,0), (1,0), (0,1)$: $\frac12\bigl((0\cdot0 - 1\cdot0) +
(1\cdot1 - 0\cdot0) + (0\cdot0 - 0\cdot1)\bigr) = \frac12$,
a \omterm{def:b2:multint:domain}{área} correta.
\end{solution}

\begin{solution}{pb:b2:multint:1}
\textbf{1.} A bola $B_n(R)$ é descrita por cotas iteradas
$-R \leq x_n \leq R$, depois $\abs{x_{n-1}} \leq \sqrt{R^2 -
x_n^2}$ e assim por diante; substituir $x_i = Ru_i$ em cada uma das
$n$ integrais a uma variável multiplica cada uma por $R$ e leva
as cotas nas de $B_n(1)$: $v_n(R) = R^n\,v_n(1) =
V_nR^n$.

\textbf{2.} Fatiando ao longo de $x_n = t$: a fatia de $B_n(1)$ é
a bola $B_{n-1}\bigl(\sqrt{1 - t^2}\bigr)$, logo, pela questão
1,
\[
V_n = \int_{-1}^1 v_{n-1}\bigl(\sqrt{1 - t^2}\bigr)\,\dd t
= V_{n-1}\int_{-1}^{1}(1 - t^2)^{\frac{n-1}2}\,\dd t .
\]

\textbf{3.} Com $t = \sin\theta$, $\dd t =
\cos\theta\,\dd\theta$ e $(1 - t^2)^{\frac{n-1}2} =
\cos^{n-1}\theta$ em $\intcc{-\pi/2}{\pi/2}$:
\[
\int_{-1}^1(1 - t^2)^{\frac{n-1}2}\dd t
= \int_{-\pi/2}^{\pi/2}\cos^n\theta\,\dd\theta
= 2\int_0^{\pi/2}\cos^n\theta\,\dd\theta = 2W_n,
\]
e $\theta \mapsto \frac\pi2 - \theta$ troca as formas em seno e
cosseno de $W_n$.

\textbf{4.} Escreva $\sin^n = \sin^{n-2}(1 - \cos^2)$ e
integre $\int\sin^{n-2}\cos\cdot\cos$ por partes ($v =
\frac{\sin^{n-1}}{n-1}$):
\[
W_n = W_{n-2} - \frac{W_n}{n-1}
\quad\Longrightarrow\quad
W_n = \frac{n-1}{n}W_{n-2}.
\]
Portanto $nW_nW_{n-1} = (n-1)W_{n-1}W_{n-2}$: a sequência
$(nW_nW_{n-1})$ é constante, igual a $1\cdot W_1W_0 =
1\cdot\frac\pi2$, logo $W_nW_{n-1} = \frac{\pi}{2n}$.

\textbf{5.} As questões 2--3 dão $V_n = 2W_nV_{n-1}$, duas vezes:
\[
V_n = 2W_n\cdot 2W_{n-1}\,V_{n-2}
= 4\,\frac{\pi}{2n}\,V_{n-2} = \frac{2\pi}n\,V_{n-2}.
\]

\textbf{6.} De $V_0 = 1$: $V_{2k} = \frac{2\pi}{2k}V_{2k-2}
= \frac\pi kV_{2k-2}$, logo $V_{2k} = \frac{\pi^k}{k!}$ por
indução. De $V_1 = 2$: $V_{2k+1} =
\frac{2\pi}{2k+1}V_{2k-1}$, logo
\[
V_{2k+1} = 2\prod_{j=1}^k\frac{2\pi}{2j+1}
= \frac{2^{k+1}\pi^k}{1\cdot3\cdot5\cdots(2k+1)} .
\]

\textbf{7.} $V_1 = 2$, $V_2 = \pi \approx 3.142$, $V_3 =
\frac{4\pi}3 \approx 4.189$, $V_4 = \frac{\pi^2}2 \approx
4.935$, $V_5 = \frac{8\pi^2}{15} \approx 5.264$, $V_6 =
\frac{\pi^3}6 \approx 5.168$, $V_7 = \frac{16\pi^3}{105}
\approx 4.725$. A razão de um passo é $V_n/V_{n-1} = 2W_n$,
e $(W_n)$ é decrescente ($\sin^n \leq \sin^{n-1}$
\omterm{def:b2:funcseq:def}{pontualmente}). Ora, $2W_5 = 2\cdot\frac45\cdot\frac23 =
\frac{16}{15} > 1$ enquanto $2W_6 =
2\cdot\frac56\cdot\frac34\cdot\frac12\cdot\frac\pi2 =
\frac{5\pi}{16} < 1$: as razões excedem $1$ até $n = 5$ e
ficam abaixo de $1$ a partir de $n = 6$ --- $(V_n)$ cresce até seu
máximo $V_5$ e depois decresce.

\textbf{8.} Para $n \geq 13 > 4\pi$: $V_n/V_{n-2} = 2\pi/n <
\tfrac12$, logo $V_{n} \leq C\cdot 2^{-n/2}$ com uma
constante fixa; melhor, para todo $q > 0$, $2\pi/n < q^2$ para $n$ grande,
de modo que $V_n/q^n \to 0$: o decaimento bate toda sequência
geométrica. A convergência de $\sum V_n$ decorre da razão
$V_n/V_{n-2} \to 0$ (compare com uma série geométrica a
partir de certa ordem).

\textbf{9.} $\sum_{k\geq0}V_{2k}x^{2k} =
\sum_{k\geq0}\frac{(\pi x^2)^k}{k!} = \eu^{\pi x^2}$, a
série exponencial (\cref{ch:b2:powerseries}), convergente para
todo $x$. Em $x = 1$: $\sum_kV_{2k} = \eu^\pi \approx
23.14$.

\textbf{10.} No cubo $\intcc{-R}R^n$ o integrando é o
produto $\prod_i\eu^{-x_i^2}$, de modo que a integral iterada
se fatora: $\bigl(\int_{-R}^R\eu^{-t^2}\dd t\bigr)^n$.
Fazendo $R \to \infty$ e usando
$\int_\R\eu^{-t^2}\dd t = \sqrt\pi$
(\cref{ex:b2:multint:polar}): $I_n = \pi^{n/2}$.

\textbf{11.} $t = u^2$ dá $\Gamma(\tfrac12) =
\int_0^\infty t^{-1/2}\eu^{-t}\dd t =
2\int_0^\infty\eu^{-u^2}\dd u = \sqrt\pi$. Iterando
$\Gamma(s+1) = s\Gamma(s)$: $\Gamma(k+1) = k!\,\Gamma(1) =
k!$, e
\[
\Gamma\Bigl(k + \frac32\Bigr)
= \Bigl(k + \frac12\Bigr)\Bigl(k - \frac12\Bigr)\cdots
\frac12\cdot\Gamma\Bigl(\frac12\Bigr)
= \frac{(2k+1)(2k-1)\cdots1}{2^{k+1}}\,\sqrt\pi .
\]

\textbf{12.} Ponha $F_n = \pi^{n/2}/\Gamma(\frac n2 + 1)$.
Como $\Gamma(\frac n2 + 1) = \frac n2\,\Gamma(\frac n2) =
\frac n2\,\Gamma(\frac{n-2}2 + 1)$, obtemos $F_n =
\frac{2\pi}nF_{n-2}$: a mesma recursão de $V_n$ (questão
5). Bases: $F_1 = \sqrt\pi/\Gamma(\frac32) =
\sqrt\pi/(\frac{\sqrt\pi}2) = 2 = V_1$ e $F_2 =
\pi/\Gamma(2) = \pi = V_2$. Por indução, $V_n = F_n$ para todo
$n$; a questão 11 transforma isso de volta nas duas \omterm{def:b2:multint:exact}{formas fechadas} da
questão 6.

\textbf{13.} Com $r = \sqrt t$: $\int_0^\infty
\eu^{-r^2}r^{n-1}\dd r = \frac12\int_0^\infty
t^{\frac n2 - 1}\eu^{-t}\dd t = \frac12\Gamma(\frac n2)$.
Portanto
\[
n\,V_n\int_0^\infty\eu^{-r^2}r^{n-1}\dd r
= V_n\cdot\frac n2\,\Gamma\Bigl(\frac n2\Bigr)
= V_n\,\Gamma\Bigl(\frac n2 + 1\Bigr) = \pi^{n/2} = I_n .
\]
Estando ambos os membros demonstrados, a identidade pode ser \emph{lida} como
a decomposição em cascas da integral gaussiana: a esfera
de raio $r$ carrega \omterm{def:b2:multint:domain}{área} $nV_nr^{n-1}$, e o peso
gaussiano $\eu^{-r^2}$ é integrado sobre as cascas.

\textbf{14.} $s_0 = V_1 = 2$ (a esfera de dimensão $0$ são dois pontos),
$s_1 = 2V_2 = 2\pi$, $s_2 = 3V_3 = 4\pi$, $s_3 = 4V_4 =
2\pi^2$; e $\int_0^R s_{n-1}r^{n-1}\dd r = V_nR^n =
v_n(R)$: a \omterm{def:b2:multint:domain}{área} é a derivada radial do volume.

\textbf{15.} Para $n = 2k$, Stirling
(\cref{thm:b2:comparison:stirling}) dá $k! \sim
\sqrt{2\pi k}\,(k/\eu)^k$, logo
\[
V_{2k} = \frac{\pi^k}{k!}
\sim \frac{(\pi\eu/k)^k}{\sqrt{2\pi k}}
= \frac1{\sqrt{\pi n}}\Bigl(\frac{2\pi\eu}n\Bigr)^{n/2}
\qquad (n = 2k).
\]
Para $n$ ímpar: $V_{2k+1} = 2W_{2k+1}V_{2k} \leq 2V_{2k}$, de modo que
valem as mesmas cotas de decaimento supergeométrico (a menos de um fator
$2$ e de um deslocamento de um no expoente) --- para todo $q >
0$, $V_n = o(q^n)$.

\textbf{16.} $V_2/4 = \pi/4 \approx 0.785$; $V_3/8 = \pi/6
\approx 0.524$; $V_{10}/2^{10} = \frac{\pi^5}{120\cdot1024}
\approx 0.0025$. Em geral $\frac{V_n/2^n}{V_{n-2}/2^{n-2}}
= \frac{2\pi}{4n} = \frac{\pi}{2n} \to 0$: a razão tende a
$0$ (supergeometricamente). A bola inscrita ocupa uma
fração que se anula: o volume do cubo migra para seus
cantos.

\textbf{17.} Pela questão 1, a bola interna de raio $1 -
\varepsilon$ tem volume $V_n(1-\varepsilon)^n$, de modo que a
casca externa carrega a fração $1 - (1 - \varepsilon)^n \to 1$.
Para $n = 100$, $\varepsilon = 0.01$: $(0.99)^{100} =
\eu^{100\ln0.99} \approx \eu^{-1.005} \approx 0.366$: cerca de
$63\%$ da bola está a menos de $1\%$ de sua superfície.

\textbf{18.} $(W_n)$ decresce, logo $W_n \leq W_{n-1} \leq
W_{n-2} = \frac{n}{n-1}W_n$: espremendo, $W_{n-1}/W_n \to 1$.
Multiplicando por $W_nW_{n-1} = \frac\pi{2n}$: $W_n^2 \sim
\frac\pi{2n}$, isto é, $W_n \sim \sqrt{\pi/(2n)}$. Cota inferior:
$W_n^2 \geq W_nW_{n+1} = \frac{\pi}{2(n+1)}$, logo $W_n \geq
\sqrt{\pi/(2(n+1))}$ para todo $n$.

\textbf{19.} Numerador: $1 - x^2 \leq \eu^{-x^2}$ dá $(1 -
x^2)^{\frac{n-1}2} \leq \eu^{-(n-1)x^2/2}$ e, com $a =
\frac{n-1}2$,
\[
\int_\delta^1(1 - x^2)^{\frac{n-1}2}\dd x
\leq \int_\delta^\infty\eu^{-ax^2}\dd x
\leq \int_\delta^\infty\frac x\delta\,\eu^{-ax^2}\dd x
= \frac{\eu^{-a\delta^2}}{2a\delta}
= \frac{\eu^{-(n-1)\delta^2/2}}{(n-1)\delta}.
\]
Denominador: $2W_n \geq \sqrt{2\pi/(n+1)}$ pela questão 18.
Dividindo obtém-se a cota exibida. Para $\delta =
s/\sqrt{n-1}$ ela se torna
$\sqrt{\tfrac{n+1}{2\pi(n-1)}}\;\eu^{-s^2/2}/s = O\bigl(
\eu^{-s^2/2}/s\bigr)$, uniformemente em $n$: fora da faixa
$\abs{x_1} \leq s/\sqrt{n-1}$ quase não há volume, para
$s$ moderadamente grande --- e, por simetria, o mesmo vale para
toda direção.

\textbf{20.} As duas afirmações coexistem porque descrevem
coordenadas diferentes do mesmo ponto. Quase todo ponto
de $B_n(1)$ tem \omterm{def:b2:nvs:norm}{norma} próxima de $1$ (questão 17: concentração
radial perto da esfera) e, no entanto, cada uma de suas $n$
coordenadas é pequena, da ordem de $1/\sqrt n$ (questão 19),
o que é coerente, pois $n$ coordenadas de tamanho $1/\sqrt n$
têm \omterm{def:b2:nvs:norm}{norma} da ordem de $1$. O volume em dimensão alta se concentra
onde todas as coordenadas repartem igualmente o orçamento de \omterm{def:b2:nvs:norm}{norma} ---
perto da esfera, mas longe de todo polo dos eixos coordenados.

\textbf{21.} Fatie $\Delta_n$ em $x_n = t \in \intcc01$: a
fatia é $\{x' \in \R^{n-1} : x_i \geq 0,\ \sum x_i \leq 1 -
t\} = (1-t)\Delta_{n-1}$, de volume
$(1-t)^{n-1}\operatorname{vol}(\Delta_{n-1})$ por homogeneidade.
Assim
\[
\operatorname{vol}(\Delta_n) =
\operatorname{vol}(\Delta_{n-1})\int_0^1(1 - t)^{n-1}\dd t =
\frac{\operatorname{vol}(\Delta_{n-1})}{n}
\quad\Longrightarrow\quad
\operatorname{vol}(\Delta_n) = \frac1{n!}\,.
\]

\textbf{22.} Os $2^n$ ortantes de sinais recortam $C_n$ em $2^n$
cópias de $\Delta_n$ (com sobreposições desprezíveis nos hiperplanos
coordenados): $\operatorname{vol}(C_n) = \frac{2^n}{n!}$. Se
$\sum\abs{x_i} \leq 1$, então $\sum x_i^2 \leq
\bigl(\sum\abs{x_i}\bigr)^2 \leq 1$: $C_n \subseteq B_n(1)$;
e $B_n(1) \subseteq \intcc{-1}1^n$ pois $\abs{x_i} \leq
\norm x$. Portanto $\frac{2^n}{n!} \leq V_n \leq 2^n$ ---
coerente com a questão 15, que situa $V_n$ entre as
escalas fatorial e geométrica.

\textbf{23.} Para $(x, y)$ no disco unitário, a fatia de
$B_4(1)$ é o disco de raio $\sqrt{1 - x^2 - y^2}$ no
plano $(z, w)$, de \omterm{def:b2:multint:domain}{área} $\pi(1 - x^2 - y^2)$. Em coordenadas
polares:
\[
V_4 = \iint_{x^2+y^2\leq1}\pi(1 - x^2 - y^2)\,\dd x\,\dd y
= \pi\int_0^{2\pi}\!\!\int_0^1(1 - \rho^2)\rho\,
\dd\rho\,\dd\alpha
= \pi\cdot2\pi\cdot\frac14 = \frac{\pi^2}2 ,
\]
em acordo com a questão 6.

\textbf{24.} A probabilidade é a razão de volumes
$\dfrac{V_{20}}{2^{20}} = \dfrac{\pi^{10}}{10!\cdot2^{20}}
\approx \dfrac{0.0258}{1\,048\,576} \approx
2.5\cdot10^{-8}$. O número de sorteios até o primeiro acerto é
da ordem do inverso, cerca de $4\cdot10^7$: um amostrador por rejeição
que funcionava lindamente para o disco ($\pi/4$ de acertos) é
inútil em dimensão $20$ --- a maldição da dimensionalidade em
uma linha.

\textbf{25.} A primeira via (Partes I--II) usou: fatiamento à Fubini
da integral iterada, substituição a uma variável
em cada coordenada (homogeneidade) e as \omterm{pb:b2:multint:1}{integrais de Wallis}
--- cálculo a uma variável puro mais indução. A segunda via
(Parte III) usou: Fubini para a estrutura de produto de $I_n$,
a mudança de variáveis polar através da integral de Gauss do
\cref{ex:b2:multint:polar} e a equação funcional da função
$\Gamma$. Elas se encontram em $V_n =
\pi^{n/2}/\Gamma(\frac n2 + 1)$, com Stirling
(\cref{thm:b2:comparison:stirling}) convertendo a fórmula
em assintótica. O volume do terceiro ano de graduação reconstrói tudo isso sobre
a integral de Lebesgue: lá Fubini e a mudança de variáveis
são teoremas para funções integráveis gerais, as coordenadas
esféricas existem em toda dimensão e os mesmos \omterm{pb:b2:multint:1}{volumes de
bolas} reaparecem como dividendos trabalhados dos
problemas de medida-produto e de Stirling --- com a convergência
dominada substituindo nossos apertos artesanais.
\end{solution}
